\section{Predictions and Falsifiability}

A theory earns its keep by sticking its neck out. A picture that explains everything after the fact and forbids nothing in advance is not a physical theory. It is a story. The vibe mesh is built to be the opposite. It commits to numbers, it names structural facts that the world either has or does not have, and it points to the one cleanest place where a single measurement could end it. This section sets out what the flow predicts, states the status of each prediction plainly, and identifies its own executioner.

Throughout, the language is the locked vocabulary of the theory. The universe is a \emph{vibe mesh}, a \emph{flow} of eight atoms. The space is the \emph{mesh}, a $\{3,4,3,4\}$ weave of \emph{docks}. A dock has exactly $24$ \emph{sites}, the $24$ directions out of it, with no rest slot and no twenty-fifth. Each site carries one \emph{tone} in $\{-1,0,+1\}$. The state law is the \emph{knit}, collide then stream. The space law is the \emph{wake}, the growing edge where new docks are born. Time is the \emph{beat}. The \emph{arrow} of time and the \emph{attraction} between bound patterns are emergent, not base.

\subsection{Quantitative predictions}

The flow puts numbers on the table. These are the hardest commitments, because a number is either right or wrong and there is nowhere to hide.

\begin{table}[ht]
\centering
\begin{tabular}{@{}lll@{}}
\toprule
prediction & value at unification & observed \\
\midrule
weak mixing angle & $\sin^2\theta_W = 3/8$ & runs to $0.231$ \\
bottom to tau mass ratio & $m_b = m_\tau$ & runs to $m_b/m_\tau \approx 2.3$ (measured $2.35$) \\
charged-lepton vs down determinant & $\det M_{\ell} = \det M_{d}$ & from discrete $\operatorname{Tr} Y = 0$ over the $16$ \\
number of generations & exactly three & exactly three \\
expansion law & $a(t) = e^{Ht}$, $H = \tfrac{1}{3}\ln R \approx 0.80$ per beat & de Sitter, $\Lambda = 3H^2 \approx 1.9$ (beat units) \\
\bottomrule
\end{tabular}
\caption{The quantitative commitments of the vibe mesh and their comparison to measurement.}
\label{tab:quantitative-predictions}
\end{table}

A few of these need their status stated plainly, because a number can be derived exactly, computed by running, or merely consistent, and these are very different epistemic claims.

\begin{principle}[The weak mixing angle is derived, not fitted]
The value $\sin^2\theta_W = 3/8$ at unification is not a fit. It is the ratio $\sum T_3^2 \big/ \sum Q^2$ summed across the sixteen states of one matter family. The numerator and the denominator are integer-and-fraction sums fixed entirely by how the sites and tones organize into the family. There is no free parameter to turn.
\end{principle}

The determinant relation $\det M_{\ell} = \det M_{d}$ is likewise derived exactly. It follows from the discrete vanishing of the hypercharge trace, $\operatorname{Tr} Y = 0$, taken over the sixteen states. A relation that holds as an identity of the discrete bookkeeping cannot be tuned away. It is a structural consequence of how charge sums over a family.

The low-energy values in the third column of Table~\ref{tab:quantitative-predictions} are \emph{computed}, then compared. The point worth stressing is that the running which carries the high-scale prediction down to laboratory energy is not imported from outside. The mechanism is internal.

\begin{result}[Coarse-graining is the running]
Radial coarse-graining of the mesh \emph{is} the renormalization-group flow. The one-dimensional decimation of the discrete tone field gives the explicit recursion for the coupling,
\[
K' = \tfrac{1}{2}\,\ln\cosh 2K ,
\]
so the beta function is a theorem about the substrate, not a postulate. Only the exact Standard Model beta coefficients remain to be supplied from the emergent gauge content.
\end{result}

The expansion rate $H \approx 0.80$ is stated in \emph{beat units}, the natural clock of the flow, one tick of the knit and one growth of the wake. The absolute rate in seconds waits on the identification of the beat with a physical time, which the theory leaves open. The shape of the law, pure exponential de Sitter growth driven by the wake, is the firm part. The conversion factor to wall-clock time is the soft part.

\subsubsection{A prediction that doubles as a probe}

The weak-mixing prediction is sharper than a single number, because the running lands at different places depending on what lives at low energy.

The high-scale value is fixed at $3/8$. Running it down with a supersymmetric low-energy spectrum lands it at $0.231$, the measured value. Running it down with the bare Standard Model content instead lands it near $0.21$, which is wrong. The same is true of the clean meeting of the three gauge couplings near $10^{16}$ GeV. That clean three-way unification requires the supersymmetric content. The Standard Model alone misses it.

\begin{proposition}[The flow predicts low-energy supersymmetry]
For the high-scale value $\sin^2\theta_W = 3/8$ to run down to the measured $0.231$, and for the three gauge couplings to meet cleanly near $10^{16}$ GeV, the low-energy spectrum must be supersymmetric. The vibe mesh therefore predicts low-energy supersymmetry. The prediction is conditional and falsifiable, since a measured low-energy world with no supersymmetry and a mixing angle the bare running cannot reach would stand against it.
\end{proposition}

This turns a successful postdiction into a live wager. The mixing angle is not only reproduced. It tells us, from inside the theory, that there should be more matter to find.

\subsection{Structural predictions}

Some predictions are qualitative but sharp. They do not name a number. They name a fact about the world that is either there or not.

\begin{result}[Matter is fermionic for free]
The spin-statistics connection is geometric in the vibe mesh, not assumed. The matter solitons are knots in the tone field of odd topological charge, and such a knot changes sign under a single full turn. That sign is the defining mark of a fermion. Matter comes out fermionic from topology alone, with no separate statistical postulate.
\end{result}

\begin{result}[A finite, frame-independent top speed]
The knit streams one dock per beat. This fixes a finite maximum signal speed, the light cone $z = 1$ in lattice units, the same in every emergent frame. A built-in top speed and an emergent Lorentz symmetry are predictions, not inputs.
\end{result}

\begin{result}[Gravity is $1/r$ in the cusp]
In the flat three-dimensional cusp of the mesh the attraction between bound patterns falls off as $1/r$, the Newtonian potential, with a $1/r^2$ holographic correlator in the bulk. The everyday inverse-square force law is an emergent consequence of how the wake's radiation pressure binds matter, not a base postulate.
\end{result}

\subsection{The cleanest falsifier: cubic anisotropy}

The mesh is a discrete crystal. A crystal is not perfectly round. This is the single sharpest experimental handle the theory has, and it is worth following the logic in full, because the energy scaling is subtle and is exactly where a careless reading would go wrong.

The light cone on the cusp is built from a discrete set of directions. It is therefore very slightly faceted rather than perfectly spherical. The faceting carries a tiny \emph{cubic anisotropy}, a direction-dependent wobble in the top speed, set by the crystal axes.

The size of this wobble grows with energy, and the power of that growth is the delicate point.

\begin{proposition}[The anisotropy is quadratic in energy]
The lowest cubic-symmetric invariant that can deform the dispersion relation is degree four in the wave vector, the sum of the fourth powers of its components, $\sum_i k_i^4$. The relative anisotropy is that degree-four term divided by the isotropic degree-two term $|k|^2$. Their ratio scales as $(k a)^2$, that is, as $(E/E_{\mathrm{Planck}})^2$, with a coefficient $c \approx 1/18 \approx 0.056$ fixed by the crystal.
\end{proposition}

The lesson is that the leading anisotropy is suppressed by two powers of the energy, not one. There is no linear term, because the lowest cubic invariant is quartic and the isotropic baseline is quadratic, so the ratio is quadratic. This pushes the effect far down.

\begin{table}[ht]
\centering
\begin{tabular}{@{}ll@{}}
\toprule
energy & fractional speed anisotropy $\delta$ \\
\midrule
$1$ GeV & $\delta \approx 4 \times 10^{-40}$ \\
GZK scale & $\delta \approx 9 \times 10^{-19}$ \\
\bottomrule
\end{tabular}
\caption{Predicted fractional anisotropy of the top speed at two energies, scaling as $(E/E_{\mathrm{Planck}})^2$ with coefficient $c \approx 1/18$.}
\label{tab:anisotropy}
\end{table}

These numbers sit far below current bounds on Lorentz violation. The quadratic, $(E/M)^2$ class of limits is weak. It requires only that the suppression scale $M$ exceed about $10^{-8}$ of the Planck mass. A Planck-scale cusp anisotropy clears that bar with enormous room to spare. So the present non-observation of Lorentz violation does not threaten the theory at all. It is exactly what the discreteness predicts at this scale.

The falsifier is nonetheless sharp, and it cuts both ways.

\begin{principle}[The named executioner]
Establishing exact Lorentz invariance to arbitrarily high energy, with no cubic anisotropy at the predicted scale, would kill the discrete vibe mesh. Conversely, a lower effective cutoff than the Planck scale would push $\delta$ upward into testable range. This is the single cleanest experimental handle the theory has, and it is a genuine handle, since the prediction is definite in both shape and size.
\end{principle}

\subsection{What else would falsify it}

A theory should list the observations it cannot survive. The vibe mesh has several.

\begin{itemize}
\item A \textbf{fourth generation} of matter. The flow forces exactly three, seeded by the order-three symmetry that permutes the three eights of the dock's $24$ sites. A fourth would have no home.
\item A \textbf{clean failure} of the running $\sin^2\theta_W$ or of the bottom-to-tau mass ratio, beyond what the open beta coefficients can absorb.
\item \textbf{No low-energy supersymmetry anywhere}, combined with a measured $\sin^2\theta_W$ that the bare Standard Model running cannot reach. This is a real tension, not a comfortable one, and it is left standing as an open problem rather than papered over.
\item A demonstration that the \textbf{knit supplies no stabilizer}, so that matter solitons would collapse and matter would not form. This would overturn the measured positive stabilizer result, which is the empirical anchor of the matter sector.
\end{itemize}

\subsection{Why it matters}

The vibe mesh is not a just-so story. It makes numbers, and it names its own executioner.

\begin{result}[The theory is falsifiable on both fronts]
The flow commits to derived numbers, the weak mixing angle $3/8$, the determinant relation, exactly three generations, the de Sitter law. And it commits to a single sharp signature that would end it, the cubic anisotropy of the top speed at the Planck scale. A theory that both predicts numbers and names the measurement that could refute it is doing the work a physical theory is supposed to do.
\end{result}

In the carrying image of the theory, these are the places to listen for whether the song is really the one the universe is singing. The numbers are the melody it has to hit. The anisotropy is the one note that, if struck wrong, would tell us the whole tune was ours and not the world's.

\subsection{See also}

The companion sections set the bounds of these claims. Scope and Coverage states what is and is not asserted. Open Problems collects the unfinished frontiers, including the supersymmetry tension named above. Implications takes up what is at stake if the predictions hold.
